% 第2章第7节：本章小结。

\section{本章小结}
\label{sec:lfo-summary}

本章围绕一个问题展开：给定一个函数，怎样稳定地变换它或从中读出一个数。第~\ref{sec:lfo-operators}节说明，线性算子是否连续由输入空间、定义域、定义域上的范数和输出空间共同决定，而不由作用公式单独决定。同一个求导规则，在一致范数下无界，在计入导数的范数下有界。算子范数给出一个对全部输入有效的误差界；强算子收敛只对每个输入分别成立，截断投影$P_n\to I$说明二者不能互换。

线性泛函是标量值的算子，它在本章中承担了双重角色。作为读数，它的良定义与连续性要先单独检验：积分读取在$L^p$上总是有界的，点值读取在一致范数下有界，在积分空间中却可能无定义或无界。在Hilbert空间中，有界读数由Riesz表示定理对应到唯一的方向；在$L^p$中对应到$L^q$中的权重，在$C(K)$中对应到测度。作为工具，泛函组成的对偶空间借Hahn--Banach定理区分了任意两个元素，又在弱收敛中充当测试对象，在伴随中把输出测试拉回输入。

弱收敛只控制逐个读数。借助Baire定理和一致有界原理，它仍然控制了序列的范数上界和极限范数的下界；在Hilbert空间和反射空间中，有界序列有弱收敛子列，而$L^1$中的质量集中说明这一结论有边界。可分前对偶的对偶空间中，有界序列有弱星收敛的子列，那些在$L^1$中找不到的极限以测度的形式出现。

伴随揭示了像与核之间的正交关系：$\ker A$是输出中看不见的输入方向，$\ker A^*$的正交补只给出像的闭包。正交投影由幂等与自伴刻画，偶函数投影和条件期望是它在函数空间和概率空间中的两个实例。逆向恢复要分别检查可解、唯一与稳定：Banach空间之间的有界双射有有界逆，但像稠密不闭的单射，例如$(x_n)\mapsto(x_n/n)$，其逆在像上无界。只凭代数上的可逆性，不能判断恢复是否稳定。

第~\ref{chap:operator-spectra-equations}章将限制到能被有限秩算子按算子范数逼近的紧算子，研究它们的谱与奇异值分解，并用本章的伴随、像核关系和有界逆讨论积分方程何时可解、何时稳定。第~\ref{chap:weak-derivatives-sobolev}章用测试函数上的线性泛函定义广义导数，第~\ref{chap:rkhs}章研究点值读数连续的Hilbert空间，并用Riesz表示构造再生核。优化目标、函数拟合与正则化则在本章的空间工具之上另行展开，本章不预先讨论它们的求解方法。
